% DRAFT: evidence-checked standalone fragment, not included by main.tex.
% Numeric provenance: pi3b_diagnostic_DRAFT_sources.json; all latency estimates unqualified.
\subsection{Physical Raspberry Pi: completed replay fails the swap gate}
\label{app:seg-pi3b-diagnostic}
A separately preregistered replay on a physical Raspberry Pi 3 Model B uses
YOLOv8n-seg, the same frozen clean/perturbed tensor pairs as the laptop replay,
and full-resolution $640\times640$ masks. The board operates under a declared
undervoltage condition at 600~MHz, rather than its nominal operating point.
The four-thread primary session completes all 20 images and 720 timing trials.
However, swap becomes active during the session, violating the preregistered
swap-off condition. The session is therefore retained as diagnostic evidence
and supplies no quality-qualified attack-latency estimate.
Both pin-load checks and the registered A/A and MDE gates pass: the geometric
mean A/A ratio is approximately 0.9985, and the image-based 80\%-power MDE
on the log-ratio scale is approximately 0.0045. These checks do not override
the swap failure; no images are removed to rescue the session.

The subsequent single-thread and component sessions both fail the swap-off
precondition before any timing trial. They are not retried under the registered
infrastructure exception. Consequently, full-cohort NMS/mask decomposition and
mask-count dose measurements are unavailable for this board. Smoke measurements
are not substituted for those missing endpoints, and separately measured
medians are not subtracted to reconstruct component costs.

Candidate and retained-mask counts match the laptop values on all 20 frozen
pairs, whereas none of the 40 clean/perturbed prediction tensors matches byte
for byte. The declared PyTorch/torchvision versions differ between devices;
finite count agreement is not numerical equivalence or general transfer
validation. This single-board diagnostic, using prepared digital tensors,
does not establish deadline failure, service denial, physical-patch efficacy,
or performance at a different clock or on another Arm device.
